\documentclass[letter,10pt]{article}
\usepackage{TLCresume}
\begin{document}

\subsection{Desarrollador de Software \hfill GEOCOM Uruguay \hfill 07/2022 --- 01/2024}

Participé en el desarrollo, mantenimiento y testing de sistemas backend y web. Trabajé en proyectos de contabilidad y reportes con Spring Framework, MongoDB, Kafka y servicios REST, y realicé mantenimiento frontend con Angular y TypeScript.

Implementé pruebas automatizadas con Selenium y Cucumber en esos proyectos, y con Java 8 y Cucumber en sistemas de punto de venta, ayudando a validar flujos críticos del negocio.

También mantuve subsistemas de promociones para aplicaciones de ventas y finanzas, corrigiendo errores e incorporando funcionalidades con Java 6 y servicios SOAP.

\vspace{5mm}

\subsection{Ingeniería de Software \hfill Proyecto Universitario}

Continuación y extensión de un sistema desarrollado por otro equipo: nuevas funcionalidades y una aplicación móvil asociada desarrollada desde cero. El proyecto implicó comprender código existente, adaptarse a decisiones de arquitectura previas y diseñar soluciones orientadas al usuario final.

\vspace{5mm}

\subsection{Programación \hfill Proyectos Universitarios}

\textbf{Taller de Programación:} aplicación en Java SE 8 y Java EE con Servlets, JSP y Tomcat, más una interfaz de administración en Swing, con metodología incremental.

\textbf{C/C++:} aplicaciones orientadas a estructuras de datos, algoritmos, programación orientada a objetos y patrones de diseño.

\vspace{5mm}

\subsection{Bases de Datos, Redes y Seguridad \hfill Proyectos Universitarios}

\textbf{Bases de Datos:} diseño de bases relacionales con modelado entidad-relación, álgebra relacional, SQL, normalización, transacciones y recuperación.

\textbf{Redes:} laboratorio con Mininet que implementa una aplicación cliente-servidor y funcionalidades de routing/forwarding de un router IP en una red emulada.

\textbf{Seguridad Informática:} laboratorios de ataques man-in-the-middle, criptografía con GPG y OpenSSL, cracking de contraseñas y escalada de privilegios, firewalls con iptables, VPN con OpenVPN y OWASP WebGoat.

\vspace{5mm}

\subsection{{\href{https://github.com/FacundoBacigalupi/cipher-chat}{CipherChat} \hfill Proyecto personal}}
Prototipo educativo de mensajería cifrada de extremo a extremo. Desarrollado con Spring Boot, React, TypeScript, PostgreSQL, JWT y WebSocket/STOMP. Implementa autenticación, chats directos, múltiples dispositivos por usuario y cifrado de mensajes por dispositivo, con claves generadas en el navegador.

\vspace{5mm}

\subsection{{\href{https://github.com/FacundoBacigalupi/home-cinema}{HomeCinema} \hfill Proyecto personal}}
Aplicación personal para organizar y reproducir contenido multimedia en red local. Backend en Spring Boot, frontend en React/TypeScript y base SQLite. Incluye escaneo de biblioteca local, progreso de reproducción, integración con TMDb y reproducción mediante HLS generado con FFmpeg.

\vspace{5mm}

\subsection{{\href{https://github.com/FacundoBacigalupi/my-stock}{Mi Stock} \hfill Proyecto personal}}
Aplicación móvil de inventario para Android con Flutter. Permite registrar productos, gestionar stock, guardar fotos, escanear códigos de barras y consultar movimientos. Usa SQLite, Riverpod y Material 3, e incluye una búsqueda visual offline experimental con TensorFlow Lite.

\vspace{5mm}

\subsection{{\href{https://github.com/FacundoBacigalupi/rubkis-cube}{Rubik's Cube 3D} \hfill Proyecto personal}}
Juego 3D en Unity 6 con solver neuronal. Implementé mecánicas de input, animaciones y lógica del cubo. Entrené una red neuronal en Python (supervised learning) exportada a ONNX y ejecutada en Unity, resolviendo mezclas de hasta 25 movimientos con una tasa de éxito cercana al 100\%.

\end{document}
